\chapter{Пространственная энергия проекторных дефектов}

Предыдущая проверка получила первое поле упорядоченной фазы,
\begin{equation}
 Q(x)=R(x)-\frac{I_3}{3},
\end{equation}
и три топологических семейства: $\mathbb Z_2$-линии, целочисленные ежи и
хопфовы текстуры. Теперь необходимо выполнить различающий энергетический
тест. Наличие ненулевой группы гомотопий ещё не делает конфигурацию
локализованной частицей.

\section{Ориентационная жёсткость полноранговой фазы}

Пусть
\begin{equation}
 R=bI_3+(a-b)P,
 \qquad P=nn^T,
 \qquad n\sim-n.
 \label{eq:v6-spatial-Q-uniaxial-form}
\end{equation}
В точке сосуществования
\begin{equation}
 a=0.9121666003\ldots,
 \qquad b=0.04391669984\ldots,
 \qquad \Delta:=a-b=0.8682499005\ldots.
\end{equation}
Так как $n\cdot dn=0$, выполнено
\begin{equation}
 \Tr(dQ\,dQ)=\Delta^2\Tr(dP\,dP)
 =2\Delta^2|dn|^2.
 \label{eq:v6-spatial-Q-gradient-identity}
\end{equation}
Следовательно, квадратичная энергия
\begin{equation}
 E_2[Q]=\frac\kappa2\int_{\mathbb R^3}\Tr(\partial_iQ\,\partial_iQ)
\end{equation}
сводится к $\kappa\Delta^2\int|\nabla n|^2$. Коэффициент $\kappa$ пока не
выведен родителем, но инфракрасный тип расходимости от его значения не
зависит.

\section{Линейный сектор}

Для минимальной дисклинации вокруг прямой оси можно взять
\begin{equation}
 n(\varphi)=
 \bigl(\cos(\varphi/2),\sin(\varphi/2),0\bigr).
\end{equation}
Знак меняется после полного обхода, но проектор возвращается к себе.
Энергия на единицу длины между радиусами ядра $\xi$ и внешним масштабом
$L$ равна
\begin{equation}
 \frac{E_{\rm line}}{\ell}
 =2\pi\kappa\Delta^2\left(\frac12\right)^2
 \log\frac L\xi.
 \label{eq:v6-line-disclination-energy}
\end{equation}
При $\kappa=1$ коэффициент логарифма равен
$1.1841572040\ldots$. Регуляризация ядра не удаляет инфракрасный
логарифм. Прямая линия в бесконечном пространстве не является частицей;
замкнутая петля имеет конечную энергию при фиксированном радиусе, но без
дополнительного баланса стремится стянуться.

\section{Точечный проекторный ёж}

Для $n(x)=x/r$ выполнено
\begin{equation}
 |\nabla n|^2=\frac2{r^2}.
\end{equation}
Поэтому энергия сферической оболочки даёт
\begin{equation}
 E_{\rm hedgehog}(L)
 =8\pi\kappa\Delta^2(L-\xi).
 \label{eq:v6-global-hedgehog-linear-energy}
\end{equation}
При единичной жёсткости линейный коэффициент равен
$18.9465152641\ldots$. Переход поля к изотропному состоянию внутри ядра
устраняет локальную сингулярность, но не дальнюю линейную расходимость.

\begin{theorem}[Глобальный проекторный ёж не локализован]
Полноранговое расщепление $1+2$ сохраняет топологический класс
$\pi_2(\mathbb{RP}^2)=\mathbb Z$, но один глобальный квадратичный член
делает энергию ежа линейно расходящейся с размером пространства.
\end{theorem}

Это уточняет результат Тома V. Спинорное накрытие по-прежнему назначает двум
ориентированным подъёмам классы $+15/-15$, но топологическая метка сама не
делает их конечномассовыми объектами.

\section{Хопфов сектор}

Хопфова текстура постоянна на пространственной бесконечности, поэтому не
имеет указанной инфракрасной расходимости. Но теорема Деррика оставляет
другую проблему. При масштабировании размера $L$
\begin{equation}
 E_2(L)=AL,
 \qquad A>0,
\end{equation}
и одна квадратичная энергия уменьшается при сжатии текстуры.

Проекторная кривизна
\begin{equation}
 F_P=P(dP)^2
\end{equation}
даёт четырёхпроизводный член с масштабированием
\begin{equation}
 E_4(L)=\frac BL,
 \qquad B>0.
\end{equation}
Тогда
\begin{equation}
 E(L)=AL+\frac BL,
 \qquad
 L_*=\sqrt{\frac BA},
 \qquad
 E_*=2\sqrt{AB}.
 \label{eq:v6-hopf-derrick-balance}
\end{equation}
Это стандартный механизм Фаддеева--Скирма и границы
Вакуленко--Капитанского. См. также
\path{arXiv:hep-th/0407245}.

Проект уже построил оба математических инварианта, но Том V доказал, что
их общая положительная относительная нормировка и размерный масштаб не
следуют из одного конечного следа $M_{35}$. Поэтому хопфова текстура есть
лучший локализуемый кандидат без нового калибровочного поля, но пока
условный.

\section{Калибровочное завершение точечного сектора}

Пространственная $SO(3)$-связность может компенсировать угловой градиент
ежа. В гладком завершении Прасад--Соммерфилда энергия конечна,
\begin{equation}
 E_{\rm BPS}=\frac{4\pi v}{g},
\end{equation}
а ориентированный оператор Каллиаса имеет индекс один. Том V уже проверил
совместимость этого механизма с хопфовой линией и Real-парой.

Однако он также доказал, что $M_{35}$ сам не порождает пространственное
дифференциальное исчисление, не выбирает
$\mathfrak{so}(3)\subset\mathfrak u(35)$ и не выводит $g$, $v$ и гладкий
полуцелый дублет ядра. Поэтому BPS-монополь остаётся математическим
образцом, а не полученной частицей проекта.

\section{Энергетическая развилка материи}

Результат оставляет два разных маршрута:
\begin{enumerate}
 \item хопфова текстура без нового калибровочного поля, но с невыведенным
 отношением $E_2/E_4$;
 \item точечный ёж с правильным классом спинорного накрытия $+15/-15$, но с
 невыведенной пространственной $SO(3)$-связностью.
\end{enumerate}
Линейная дисклинация может быть струной или границей домена, но не
элементарной точечной частицей в $3+1$ измерениях.

Следующая проверка должна проверить более жёсткий и топологически прямой
маршрут: можно ли получить калибровочное завершение ежа из уже
существующей пространственно-внутренней суперсвязности и данных спинорного
накрытия, не вводя вручную новую группу и новый дублет.

\begin{nogocriterion}[Три запрета преждевременной частицы]
Линия с инфракрасно расходящейся энергией, глобальный ёж с линейной
расходимостью и хопфова текстура без фиксированного масштаба не могут быть
объявлены физическими частицами.
\end{nogocriterion}

\begin{protocol}
\begin{enumerate}
 \item $\mathbb Z_2$-линия: \textbf{логарифмически нелокальна};
 \item точечный глобальный ёж: \textbf{линейно нелокален};
 \item классы спинорного накрытия ежа: \textbf{$+15/-15$ сохраняются};
 \item хопфова текстура при одном $E_2$: \textbf{коллапсирует};
 \item хопфова текстура при $E_2+E_4$: \textbf{имеет конечный радиус};
 \item отношение $E_2/E_4$ из родителя: \textbf{не выведено};
 \item гладкое калибровочное завершение ежа: \textbf{математически есть};
 \item пространственная $SO(3)$-динамика проекта: \textbf{не выведена};
 \item полученная наблюдаемая частица: \textbf{нет};
 \item следующая проверка:
 \textbf{родительское калибровочное завершение спинорного накрытия};
 \item рождение материи: \textbf{не закрыто}.
\end{enumerate}
\end{protocol}

Машинный сертификат сохранён в
\path{s2t_v6_spatial_projective_defect_energy_spectrum_gate_results.json}.